% 图契约：串联线圈的助磁与反磁；图中的端口、几何、箭头和方程对应本题。
% 版面：沿用前三模块的几何图+推导分区；目标印宽17cm，简图至中等密度。
% 布局：左侧真实结构，右侧关键关系；真实电路边不作为架构流程箭头处理。
\documentclass[tikz,border=7pt]{standalone}
\usepackage{fontspec,xeCJK,amsmath,amssymb}
\IfFontExistsTF{Times New Roman}{\setmainfont{Times New Roman}}{\setmainfont{TeX Gyre Termes}}
\IfFontExistsTF{Heiti SC}{\setCJKmainfont[AutoFakeBold=1.5]{Heiti SC}}{\setCJKmainfont[AutoFakeBold=1.5]{FandolHei-Regular}}
\usetikzlibrary{arrows.meta,bending,calc,patterns,positioning,decorations.pathmorphing}
\definecolor{ink}{HTML}{243745}
\definecolor{blue}{HTML}{275DAB}
\definecolor{green}{HTML}{177C63}
\definecolor{red}{HTML}{BD4935}
\definecolor{orange}{HTML}{AA6C16}
\definecolor{purple}{HTML}{715393}
\tikzset{every node/.style={font=\fontsize{9.3}{12.5}\selectfont,text=ink},
arrow/.style={-{Stealth[length=5pt 1.5,width'=0pt 3]},shorten >=2pt,shorten <=1pt,line width=1pt},
arrow short/.style={-{Stealth[length=3pt 1.0,width'=0pt 2.5]},shorten >=1pt,shorten <=0pt,line width=.8pt},
axis/.style={arrow,draw=ink!60,line width=.6pt},
vec/.style={arrow,draw=blue},
wire/.style={draw=ink!75,line width=.8pt},
guide/.style={dashed,draw=ink!45,line width=.5pt},
card/.style={rounded corners=2pt,draw=blue!65,fill=blue!4,align=center,text width=6.5cm,inner xsep=9pt,inner ysep=7pt},
title/.style={font=\fontsize{10.5}{14}\selectfont\bfseries,text width=15cm,align=center},
label/.style={font=\fontsize{8.4}{11}\selectfont,align=center},
note/.style={label,text width=6.1cm}}
\begin{document}
\begin{tikzpicture}[x=1cm,y=1cm]
\path[use as bounding box] (-8,-3.5) rectangle(8,3.65);
\node[title] at (0,3.25) {理论图\quad 串联线圈的助磁与反磁};
\begin{scope}[shift={(-4.05,0)}]\begin{scope}[shift={(0,1.3)}]\draw[wire](-2.8,0)--(-2.5,0);\draw[wire](-2.5,0)..controls(-2.5,0.5)and(-2.025,0.5)..(-2.025,0);\draw[wire](-2.025,0)..controls(-2.025,0.5)and(-1.55,0.5)..(-1.55,0);\draw[wire](-1.55,0)..controls(-1.55,0.5)and(-1.0750000000000002,0.5)..(-1.0750000000000002,0);\draw[wire](-1.0750000000000002,0)..controls(-1.0750000000000002,0.5)and(-0.6000000000000001,0.5)..(-0.6000000000000001,0);\draw[wire](-0.6,0)--(-0.3,0);\node[below]at(-1.5499999999999998,-0.2){$L_1$};\draw[wire](0.3,0)--(0.6,0);\draw[wire](0.6,0)..controls(0.6,0.5)and(1.075,0.5)..(1.075,0);\draw[wire](1.075,0)..controls(1.075,0.5)and(1.5499999999999998,0.5)..(1.5499999999999998,0);\draw[wire](1.5499999999999998,0)..controls(1.5499999999999998,0.5)and(2.025,0.5)..(2.025,0);\draw[wire](2.025,0)..controls(2.025,0.5)and(2.5,0.5)..(2.5,0);\draw[wire](2.5,0)--(2.8,0);\node[below]at(1.5499999999999998,-0.2){$L_2$};\draw[wire](-.3,0)--(.3,0);\fill[orange](-2.5,.65)circle(2.5pt);\fill[orange](0.6,.65)circle(2.5pt);\node[above]at(0,1){助磁：电流向右};\end{scope}\begin{scope}[shift={(0,-1)}]\draw[wire](-2.8,0)--(-2.5,0);\draw[wire](-2.5,0)..controls(-2.5,0.5)and(-2.025,0.5)..(-2.025,0);\draw[wire](-2.025,0)..controls(-2.025,0.5)and(-1.55,0.5)..(-1.55,0);\draw[wire](-1.55,0)..controls(-1.55,0.5)and(-1.0750000000000002,0.5)..(-1.0750000000000002,0);\draw[wire](-1.0750000000000002,0)..controls(-1.0750000000000002,0.5)and(-0.6000000000000001,0.5)..(-0.6000000000000001,0);\draw[wire](-0.6,0)--(-0.3,0);\node[below]at(-1.5499999999999998,-0.2){$L_1$};\draw[wire](0.3,0)--(0.6,0);\draw[wire](0.6,0)..controls(0.6,0.5)and(1.075,0.5)..(1.075,0);\draw[wire](1.075,0)..controls(1.075,0.5)and(1.5499999999999998,0.5)..(1.5499999999999998,0);\draw[wire](1.5499999999999998,0)..controls(1.5499999999999998,0.5)and(2.025,0.5)..(2.025,0);\draw[wire](2.025,0)..controls(2.025,0.5)and(2.5,0.5)..(2.5,0);\draw[wire](2.5,0)--(2.8,0);\node[below]at(1.5499999999999998,-0.2){$L_2$};\draw[wire](-.3,0)--(.3,0);\fill[orange](-2.5,.65)circle(2.5pt);\fill[orange](2.5,.65)circle(2.5pt);\node[above]at(0,1){反磁：电流向右};\end{scope}\end{scope}\node[card,draw=blue!65,fill=blue!4]at(3.9,1.8){$L_{\rm aid}=L_1+L_2+2M$};\node[card,draw=green!65,fill=green!4]at(3.9,0){$L_{\rm opp}=L_1+L_2-2M$};\node[card,draw=orange!65,fill=orange!4]at(3.9,-1.8){橙点指明端口相对极性，$M\ge0$};
\end{tikzpicture}\end{document}